\section{Background and Related Work}\label{background-and-related-work}

\subsection{Biomedical Background of the Task}\label{biomedical-background-of-the-task}

Human induced pluripotent stem cell-derived cardiomyocytes (hiPSC-CMs) are widely used as an in vitro model for cardiac biology, disease modeling, and phenotypic screening \cite{ahmed2020hipscm_maturation}. In these applications, the morphology of individual cells is not a cosmetic detail: it is closely tied to questions about maturation state, structural organization, drug response, and culture quality. This is one of the reasons why a cell-level segmentation task is meaningful in the first place. The goal is not merely to detect whether cardiomyocytes are present, but to recover individual cell extents in a way that supports downstream quantitative analysis.

Compared with many standard microscopy benchmarks, cardiomyocytes are structurally awkward segmentation targets. They are often elongated, anisotropic, irregular in width, and densely packed. Neighboring cells can touch or partially overlap in projection, while local boundaries may be weak in brightfield imaging. As a result, the object of interest is visually more complex than round or nucleus-like cells, and the mismatch between ``where a cell is'' and ``what its full outline should be'' is larger than in many easier segmentation settings.

This scope also sets an important boundary for the thesis. The current dataset and all locked results concern \textbf{human hiPSC-derived cardiomyocytes}. Therefore, the claims made in this thesis should be read as human hiPSC-CM claims. Possible transfer to mouse cardiomyocytes is scientifically interesting, but it remains a future adaptation question rather than a present conclusion of this work.

\subsection{Human--Mouse Transfer Boundary and Species-Specific Differences}\label{human-mouse-transfer-boundary-and-species-specific-differences}

The practical transfer question is not whether human-to-mouse adaptation is imaginable, but whether current human hiPSC-CM evidence can be extrapolated without species-aware adjustment. Existing evidence suggests this is unsafe: rodents and humans differ in electrophysiology and ion-channel dominance \cite{joukar2021rodent_human_ecg}, while hiPSC-CMs are typically immature relative to adult cardiomyocytes \cite{ahmed2020hipscm_maturation}. Cardiac single-cell and spatial-omics reviews further emphasize species-dependent cellular composition and cardiomyocyte state differences that affect interpretation across species \cite{palmer2024cardiac_omics}. A detailed species-gap summary is provided in Appendix~\ref{scope-and-transfer-notes}.

Therefore, the thesis treats mouse transfer as a future adaptation track rather than a supported claim of direct zero-shot reuse. Any credible transfer route should include at least: (i) mouse-domain zero-shot audit, (ii) detector-prior retuning, and (iii) mouse-specific decoder fine-tuning under the same split-aware evaluation protocol.

\subsection{Imaging Modalities and Structural Cues}\label{imaging-modalities-and-structural-cues}

The Allen hiPSC-CM dataset used in this thesis combines several channels with different biological meanings. Brightfield provides the main low-cost imaging modality and is central to the deployment-oriented question of this thesis, but it also gives the weakest boundary information. The alpha-Actinin2 channel highlights sarcomeric or cytoskeletal structure inside cardiomyocytes, which is highly informative for cell identity and internal organization but does not directly provide a clean whole-cell contour. The DAPI channel highlights nuclei and is therefore useful for localization, yet a nucleus is only a proxy for a whole cardiomyocyte and not its full spatial extent.

These channels are therefore complementary rather than interchangeable. Brightfield is attractive because it is broadly available and less burdensome than fluorescence-heavy acquisition. Actn2 provides rich structural information about cardiomyocyte organization. DAPI provides a strong positional prior for candidate cells. The central challenge is that none of these signals alone perfectly defines the whole-cell boundary, which is why both detection quality and segmentation quality must be considered jointly in this project.

From the perspective of method design, these channel semantics are not merely descriptive background. They define the biological priors that motivate the pipeline itself. DAPI provides a localization prior for candidate-cell prompting, while Actn2 provides a cardiomyocyte-specific structural prior that helps distinguish true cardiomyocyte extent from generic brightfield texture. This is why the thesis treats channel handling and prompt construction as part of the method rather than as minor preprocessing details.

\subsection{Why Whole-Cell Instance Segmentation Is Difficult in Cardiomyocytes}\label{why-whole-cell-instance-segmentation-is-difficult-in-cardiomyocytes}

The present thesis focuses on whole-cell instance segmentation rather than nucleus segmentation or semantic foreground extraction. This choice is important biologically and methodologically. A nucleus-level mask is insufficient for morphology analysis, and a semantic foreground mask can overestimate performance by merging adjacent cells into a single blob. For cardiomyocytes, this failure mode is particularly plausible because elongated neighboring cells can align, touch, or overlap in ways that make the global foreground look reasonable while the instance structure is wrong.

This difficulty also explains several design choices in the project. DAPI-based detection can provide candidate boxes, but nuclear localization does not resolve the full cell boundary. Structural channels such as Actn2 can provide clues about internal organization, but they do not automatically separate touching cells. Brightfield imaging is the most practical modality for future deployment, but its weak boundaries make direct whole-cell delineation difficult. The thesis therefore studies a prompt-based segmentation pipeline in which localization and per-cell mask prediction are treated as related but distinct problems.

\subsection{Foundation Models for Interactive Segmentation}\label{foundation-models-for-interactive-segmentation}

Segment Anything Model (SAM) established a new paradigm for promptable segmentation by combining a Vision Transformer image encoder, a prompt encoder, and a lightweight mask decoder \cite{kirillov2023sam}. Its main contribution was not a single-task architecture, but a reusable segmentation representation learned from large-scale pretraining. For biomedical imaging, this idea is attractive because annotation is expensive, morphology is diverse, and many applications require transferring a model across domains with limited labeled data.

However, direct transfer from natural-image pretraining to microscopy is not guaranteed to work well. Microscopy images differ in texture, intensity statistics, scale, and object structure. In particular, cardiomyocytes are long, anisotropic, and densely arranged, which violates many of the geometric assumptions that make generic segmentation models perform well on more regular biological objects. This motivates domain-specific adaptation rather than direct reuse of a released checkpoint.

\subsection{CellSAM and the Problem of Released-Checkpoint Semantics}\label{cellsam-and-the-problem-of-released-checkpoint-semantics}

CellSAM extends SAM to the cell segmentation domain by combining a detection branch and a segmentation branch \cite{israel2024cellsam}. Conceptually, its published methodology describes a two-stage pipeline: first learning cell detection, then adapting the segmentation branch. In practice, however, a released checkpoint is not only a scientific object but also a software object. The exact meaning of ``shared backbone'', ``stage 2 branch'', or released inference branch semantics must be verified against the released implementation.

The published CellSAM paper is specific about the intended Stage 2 training role: after detection training, the ViT and the SAM mask decoder are fixed, and only the neck is fine-tuned under ground-truth box and segmentation supervision. This published description is important, but it does not fully determine the semantics of the released checkpoint. The public weights still need to be audited at code level and parameter level before they can be used as a thesis baseline.

This distinction matters for the present thesis. Our project-level audit shows that the public CellSAM checkpoint contains multiple SAM branches with different parameter sets and roles. Therefore, cardiomyocyte adaptation cannot be described simply as ``fine-tuning CellSAM'' without specifying which branch is used for inference, which branch is aligned with CellFinder, and which branch is actually being optimized during our experiments. This is one of the reasons the thesis treats model-path auditing as part of the method rather than a minor implementation note.

\subsection{Medical-Domain Adaptation of SAM}\label{medical-domain-adaptation-of-sam}

Medical-domain variants of SAM illustrate different strategies for adapting a foundation model beyond its original training domain. MedSAM emphasizes large-scale medical pretraining and shows that SAM-style architectures can become strong medical segmentation baselines when trained on sufficiently broad medical data \cite{ma2024medsam}. This makes MedSAM a useful upper-bound comparison in our work: it represents what can be achieved when domain transfer is supported by much larger medical pretraining resources.

Another line of work focuses on parameter-efficient adaptation rather than full-domain pretraining. Techniques such as LoRA and adapter-based tuning modify only a small subset of parameters while keeping most of the backbone fixed \cite{hu2022lora}. These approaches are attractive in small-data regimes because they reduce computational cost and can limit overfitting. In our project, LoRA is explored as a historical side path, but it is not the current thesis mainline because the sealed segmentation route currently uses decoder-focused T28 under dual-seed validation reporting.

\subsection{Model Architecture and Algorithmic Principles for This Thesis}\label{model-architecture-and-algorithmic-principles-for-this-thesis}

The final thesis pipeline is a cascaded detector-to-segmenter system, with explicit role separation between box generation and mask prediction. This separation is intentional: it allows us to attribute errors to prompt quality versus mask refinement quality.

\textbf{CellSAM foundation architecture.}
CellSAM follows the SAM family design \cite{kirillov2023sam,israel2024cellsam}: a ViT image encoder produces image embeddings, a prompt encoder embeds spatial prompts (boxes or points), and a mask decoder predicts instance masks. In the released checkpoint used by this thesis, segmentation inference is run from the released \texttt{model\_cp} branch (single-mask decoding), and this branch is adapted by the T28 decoder-focused fine-tuning route.

\textbf{Detector--segmenter role split in this project.}
In the final deployment line, detection prompts are generated by \texttt{T33g}, which is the CellSAM CellFinder detector branch (AnchorDETR-style transformer detection head) under candidate-aware DAPI+Actn2 priors \cite{israel2024cellsam}. Per-box mask refinement is produced by \texttt{T28}. Therefore, the detector and segmenter are structurally related but operationally decoupled in training and inference. This is why end-to-end results are interpreted as a prompt-quality and segmentation-quality joint outcome rather than as a single-module score.

\textbf{Cellpose baseline principle (v4.0.1 in this thesis).}
Cellpose is a strong generalist cell-segmentation baseline built around flow-based mask reconstruction and broad cross-cell-type transfer behavior \cite{stringer2021cellpose}. In this thesis it serves as a transfer stress test: whether a robust generic cell segmenter can directly handle dense hiPSC-cardiomyocyte morphology without cardiomyocyte-specific prompt design.

\textbf{MedSAM baseline principle.}
MedSAM keeps the SAM-style promptable segmentation structure but relies on large-scale medical-domain pretraining \cite{ma2024medsam}. Its role here is not to mimic our pipeline design, but to provide a strong external foundation baseline. The contrast is methodological: MedSAM emphasizes broad medical pretraining scale, whereas our line emphasizes released-checkpoint audit, biology-prior prompts, and cardiomyocyte-specific adaptation under constrained in-domain data.

\subsection{Classical and Application-Oriented Cell Segmentation Baselines}\label{classical-and-application-oriented-cell-segmentation-baselines}

The thesis also positions the proposed approach against established cell segmentation baselines. Cellpose is important because it is a widely adopted generalist cell segmentation method and also appears in the CellSAM evaluation ecosystem \cite{stringer2021cellpose}. For this reason, replacing the project's earlier mismatched Cellpose route with a version-controlled rerun is necessary for a fair comparison. In the current thesis evidence tier, the T31 rerun serves exactly this role.

Beyond end-to-end segmentation models, classical cell-centric analysis workflows still have practical value for quality control and audit. In particular, CellProfiler-style pipelines provide reusable primitives for nucleus/cell association, object-level morphology statistics, and feature-based quality control \cite{carpenter2006cellprofiler,mcquin2018cellprofiler3}. In this thesis, this family is not treated as a direct competing end-to-end baseline; instead, its value is methodological support for annotation audit and detector-candidate screening rules (for example, identifying suspicious nucleus-to-cell mismatches and likely non-cardiomyocyte outliers).

The thesis does not need every baseline ever run in the project to make its argument. It needs a set of comparisons that are methodologically clean, auditable, and informative for the cardiomyocyte setting. In this thesis, the baseline roles are deliberately non-overlapping: Cellpose tests direct transfer of a strong generalist cell segmenter; SAM ViT-B tests whether generic promptable segmentation plus GT boxes is already sufficient; CellSAM-pretrained tests zero-shot transfer from a released cell-domain foundation model; and MedSAM tests whether broad medical-domain pretraining dominates a cardiomyocyte-specific route. The role of the audited adaptation line (T28 mainline, with T27a retained as a locked historical reference) is different from all four: it asks whether a domain-audited cardiomyocyte-specific path can exceed these comparison axes while remaining split-aware and reproducible.

\subsection{Historical Cardiomyocyte Analysis and Segmentation Methods}\label{historical-cardiomyocyte-analysis-and-segmentation-methods}

Direct precedents for the exact task studied in this thesis are surprisingly sparse. Much of the earlier cardiomyocyte-image literature does not solve 2D whole-cell instance segmentation of dense hiPSC-derived cardiomyocyte cultures with automatic prompts. Instead, the literature is distributed across several partially related problem classes.

First, a substantial body of work focuses on \textbf{cardiomyocyte nuclei identification or classification} rather than whole-cell extent recovery. A representative example is Nuquantus, which applies machine learning to identify and characterize nuclei in complex cardiac immunofluorescence tissue images \cite{gross2016nuquantus}. Such systems are biologically useful, but they solve a different task from whole-cell instance segmentation because nuclei provide localization priors, not complete cell boundaries.

Second, some works provide \textbf{semi-automated quantitative histology pipelines} for cardiac tissue, such as HeartJ-style analysis \cite{droste2023heartj}. These pipelines are valuable for measuring cell size, fibrosis, or tissue organization, but they typically depend on substantial preprocessing, user intervention, or application-specific image conditions. They should therefore be understood as cardiac image-analysis tools rather than direct counterparts to the present prompt-based whole-cell segmentation pipeline.

Third, there are \textbf{application-oriented deep-learning methods} for microscopy cell segmentation that can be applied to cardiac images, but they are not cardiomyocyte-specialist foundation models. Cellpose is the clearest example in this category: it is a strong generalist cell segmenter and therefore a meaningful baseline, yet it is not designed specifically around cardiomyocyte morphology or around the prompt-based detector-plus-segmenter setting used in CellSAM-derived pipelines.

Fourth, some recent work is much closer biologically but still differs technically from the present task. For example, recent ventricular-cardiomyocyte segmentation toolkits target 3D confocal or volumetric settings. These methods are relevant because they show that cardiomyocyte segmentation is an active research problem, but they address a different imaging regime from the 2D brightfield-plus-fluorescence hiPSC-CM setting of this thesis.

Taken together, this literature gap strengthens the methodological motivation of the thesis. The main challenge is not simply to outperform a weak historical baseline. It is to adapt a released foundation-style cell segmentation artifact to a domain in which (i) nuclei are informative but incomplete, (ii) structural channels encode cardiomyocyte identity without giving clean whole-cell contours, and (iii) deployment depends on prompt quality as much as on mask quality.

\subsection{Dataset-Source Boundary}\label{dataset-source-boundary}

The Allen ecosystem is also important context, but it must be interpreted carefully \cite{rafelski2024allen}. Allen Cell Feature Explorer and related Allen Cell resources expose segmented cell data, image-derived features, and structure-focused analysis pipelines for hiPSC and hiPSC-derived cardiomyocyte data. Allen also publishes the Allen Cell Structure Segmenter, which is a toolkit for 3D intracellular and cellular structure segmentation and supports both image-processing pipelines and machine-learned models for selected structures.

However, this is not the same problem as the one studied here. The Allen public tooling primarily addresses 3D cell-structure segmentation and cell-state exploration, whereas the present thesis studies 2D whole-cardiomyocyte instance segmentation from brightfield, DAPI, and Actn2 channels under a detector-to-segmentation workflow. This distinction matters because it prevents overclaiming that the data source already provides a solved whole-cell segmentation route for the exact deployment setting considered in this project.

\subsection{Summary of the Gap Addressed in This Thesis}\label{summary-of-the-gap-addressed-in-this-thesis}

The literature suggests three unresolved tensions that motivate this work.

First, foundation-model transfer is promising but not automatically reliable for cardiomyocytes. Second, strong performance claims require careful separation between model quality, prompt quality, and inference-path choice. Third, released checkpoints may encode more implementation complexity than the corresponding paper narrative makes obvious.

This thesis addresses these tensions by combining a code-level audit of the released CellSAM artifact with a cardiomyocyte-specific fine-tuning mainline, biology-prior-guided input and prompting design, and a split-aware evaluation strategy. The result is not only a stronger segmentation model, but also a cleaner statement of what is actually supported by the current evidence. Under this framing, the thesis contribution is stronger than a simple ``better PQ'' claim: it is about artifact auditing, domain-specific adaptation, biology-prior-driven design, and evidence-tiered evaluation on a difficult human hiPSC-CM task.
